% GENERATED FILE. Edit canonical lessons, then run npm run generate:books.
% canonical-source-hash: fnv1a64:0a1ee0ad565ca67a
% canonical-lessons: BN-C253-mitthe, BN-C253-rohosyo, BN-C253-rosikota, BN-C253-hasi, BN-C253-bondhutvo

\chapter{Lie, Mystery, Joke, Smile, Friendship}
\label{ch:bn-lie-mystery-joke-smile-friendship}
\hlchaptermodality{253}

\begin{chapteropening}
\textbf{By the end of this chapter:} \emph{I can say a lie, a mystery, a joke, a smile and friendship.}

Complete the last lesson of chapter 253: I can say a lie, a mystery, a joke, a smile and friendship.
\end{chapteropening}

\section[\texorpdfstring{\bn{মিথ্যে} (mitthe)}{মিথ্যে (mitthe)}]{\bn{মিথ্যে} (mitthe) — a lie}

\label{lesson:BN-C253-mitthe}



\begin{quote}
Before the new one: say the Bengali for a granddaughter, then the Bengali for a father-in-law.
\end{quote}

\subsection*{You'll want to know: \bn{মিথ্যে}}
\textbf{\bn{মিথ্যে}} — \emph{mitthe} — \textquotedblleft{}a lie\textquotedblright{}.

\begin{grammarlens}[title={What you've built}]
One more word for this chapter.
\end{grammarlens}

\subsection*{Your turn}
\begin{itemize}
\raggedright
  \item \emph{Say it:} \emph{mitthe}
  \item \emph{Say it:} \emph{mitthe}, once more
  \item \emph{Say it:} say \emph{mitthe}
  \item \emph{Recall:} say the Bengali for a granddaughter, then the Bengali for a father-in-law, then say \emph{mitthe} again
\end{itemize}

\begin{tcolorbox}[breakable,colback=teal!4,colframe=teal!35!black,title={Before you move on}]
What does \bn{মিথ্যে} mean? (\textquotedblleft{}A lie\textquotedblright{}.) Say it once more.
\end{tcolorbox}
\section[\texorpdfstring{\bn{রহস্য} (rôhôsyô)}{রহস্য (rôhôsyô)}]{\bn{রহস্য} (rôhôsyô) — a mystery}

\label{lesson:BN-C253-rohosyo}



\begin{quote}
Before the new one: say the Bengali for a father-in-law, then the Bengali for a lie.
\end{quote}

\subsection*{You'll want to know: \bn{রহস্য}}
\textbf{\bn{রহস্য}} — \emph{rôhôsyô} — \textquotedblleft{}a mystery\textquotedblright{}.

\begin{grammarlens}[title={What you've built}]
One more word for this chapter.
\end{grammarlens}

\subsection*{Your turn}
\begin{itemize}
\raggedright
  \item \emph{Say it:} \emph{rôhôsyô}
  \item \emph{Say it:} \emph{rôhôsyô}, once more
  \item \emph{Say it:} say \emph{rôhôsyô}
  \item \emph{Recall:} say the Bengali for a father-in-law, then the Bengali for a lie, then say \emph{rôhôsyô} again
\end{itemize}

\begin{tcolorbox}[breakable,colback=teal!4,colframe=teal!35!black,title={Before you move on}]
What does \bn{রহস্য} mean? (\textquotedblleft{}A mystery\textquotedblright{}.) Say it once more.
\end{tcolorbox}
\section[\texorpdfstring{\bn{রসিকতা} (rôsikôtā)}{রসিকতা (rôsikôtā)}]{\bn{রসিকতা} (rôsikôtā) — a joke}

\label{lesson:BN-C253-rosikota}



\begin{quote}
Before the new one: say the Bengali for a lie, then the Bengali for a mystery.
\end{quote}

\subsection*{You'll want to know: \bn{রসিকতা}}
\textbf{\bn{রসিকতা}} — \emph{rôsikôtā} — \textquotedblleft{}a joke\textquotedblright{}.

\begin{grammarlens}[title={What you've built}]
One more word for this chapter.
\end{grammarlens}

\subsection*{Your turn}
\begin{itemize}
\raggedright
  \item \emph{Say it:} \emph{rôsikôtā}
  \item \emph{Say it:} \emph{rôsikôtā}, once more
  \item \emph{Say it:} say \emph{rôsikôtā}
  \item \emph{Recall:} say the Bengali for a lie, then the Bengali for a mystery, then say \emph{rôsikôtā} again
\end{itemize}

\begin{tcolorbox}[breakable,colback=teal!4,colframe=teal!35!black,title={Before you move on}]
What does \bn{রসিকতা} mean? (\textquotedblleft{}A joke\textquotedblright{}.) Say it once more.
\end{tcolorbox}
\section[\texorpdfstring{\bn{হাসি} (hāsi)}{হাসি (hāsi)}]{\bn{হাসি} (hāsi) — a smile, laughter}

\label{lesson:BN-C253-hasi}



\begin{quote}
Before the new one: say the Bengali for a mystery, then the Bengali for a joke.
\end{quote}

\subsection*{You'll want to know: \bn{হাসি}}
\textbf{\bn{হাসি}} — \emph{hāsi} — \textquotedblleft{}a smile, laughter\textquotedblright{}.

\begin{grammarlens}[title={What you've built}]
One more word for this chapter.
\end{grammarlens}

\subsection*{Your turn}
\begin{itemize}
\raggedright
  \item \emph{Say it:} \emph{hāsi}
  \item \emph{Say it:} \emph{hāsi}, once more
  \item \emph{Say it:} say \emph{hāsi}
  \item \emph{Recall:} say the Bengali for a mystery, then the Bengali for a joke, then say \emph{hāsi} again
\end{itemize}

\begin{tcolorbox}[breakable,colback=teal!4,colframe=teal!35!black,title={Before you move on}]
What does \bn{হাসি} mean? (\textquotedblleft{}A smile, laughter\textquotedblright{}.) Say it once more.
\end{tcolorbox}
\section[\texorpdfstring{\bn{বন্ধুত্ব} (bôndhutvô)}{বন্ধুত্ব (bôndhutvô)}]{\bn{বন্ধুত্ব} (bôndhutvô) — friendship}

\label{lesson:BN-C253-bondhutvo}



\begin{quote}
Before the new one: say the Bengali for a joke, then the Bengali for a smile.
\end{quote}

\subsection*{You'll want to know: \bn{বন্ধুত্ব}}
\textbf{\bn{বন্ধুত্ব}} — \emph{bôndhutvô} — \textquotedblleft{}friendship\textquotedblright{}.

\begin{grammarlens}[title={What you've built}]
One more word for this chapter.
\end{grammarlens}

\subsection*{Your turn}
\begin{itemize}
\raggedright
  \item \emph{Say it:} \emph{bôndhutvô}
  \item \emph{Say it:} \emph{bôndhutvô}, once more
  \item \emph{Say it:} say \emph{bôndhutvô}
  \item \emph{Recall:} say the Bengali for a joke, then the Bengali for a smile, then say \emph{bôndhutvô} again
\end{itemize}

\begin{tcolorbox}[breakable,colback=teal!4,colframe=teal!35!black,title={Before you move on}]
What does \bn{বন্ধুত্ব} mean? (\textquotedblleft{}Friendship\textquotedblright{}.) Say it once more.
\end{tcolorbox}
